\documentclass[10pt,a4paper]{amsart}
\usepackage[margin=19mm]{geometry}
\usepackage{amsmath,amssymb,mathtools}
\usepackage[scr=rsfs,cal=euler]{mathalfa}
\usepackage{xfrac}
\usepackage[unicode,hidelinks,bookmarks=false]{hyperref}
\newcommand{\ie}{i.e.\ }
\newcommand{\cf}{cf.\ }
\newcommand{\resp}{respectively\ }
\newcommand{\Subsection}{\subsection}
\providecommand{\nobreakdash}{\nobreak\hbox{-}}
\setlength{\emergencystretch}{2em}
\multlinegap0pt
\usepackage{tikz}
\usepackage{amssymb}
\usepackage{float}
\newtheorem{proposition}{Proposition}
\title{S-packing chromatic critical paths and cycles}
\author{G\"ulnaz Boruzanl{\i} Ekinci, Csilla Bujt\'as, Didem G\"oz\"upek, Asl{\i}han G\"ur}
\date{Source: arXiv:2604.02267v1 (CC BY 4.0)}
\begin{document}
\maketitle

\noindent\emph{Source and licence.} S-packing chromatic critical paths and cycles. Version arXiv:2604.02267v1 (CC BY 4.0), distributed by arXiv under the Creative Commons Attribution 4.0 International licence, http://creativecommons.org/licenses/by/4.0/, as recorded in the arXiv licence metadata for this exact version and shown on the abstract page https://arxiv.org/abs/2604.02267v1. Full licence notice: \texttt{licenses/arxiv-2604.02267-CC-BY-4.0.txt}.\par\smallskip

\noindent\textit{Editorial scope.} Counted unit: the article's proposition proposition (source0.tex lines 1035--1040). The article's complete original proof of it is delivered with it (source0.tex lines 1042--1095, one \texttt{proof} environment of 54 lines). No mathematics is written, repaired or re-derived: the statement and the proof are the article's own lines, in the article's own notation, and no retained line is substituted or re-flowed.  No further source-local context is needed: every cross-reference of every retained line resolves inside the two delivered spans.  Nothing was omitted from any retained span, so the package carries no omission record.  No retained line contains a citation, so the wrapper carries no bibliography and no bibliography entry.  No retained line contains a figure environment, an image-inclusion command or drawing code, so no image is delivered and the package declares no asset dependency.  Editorial formatting only, in addition: an AMS \texttt{amsart} preamble that restates the article's own declarations and macros used by the retained lines, namely the environments \texttt{proposition} on separate counters, together with the standard packages those lines need.  The retained environments print in this document as Proposition 1 in order of appearance, whereas the article prints the same items with the numbering of its own full document.  The complete native source of the article as distributed in the arXiv e-print for this exact version is bundled unchanged as \texttt{sources/197724/source0.tex}.
\begin{proposition}
	If $S \in \mathcal{S}_{1,2,4,6}$, then $C_n$ is $S$-packing critical if and only if $n \in \{3, 5, 6, 7, 9, 10, 11,\allowbreak 12, 13, 17, 18, 19, 20, 25, 26, 27, 33, 34, 41\}$. 
	
	Equivalently, if $S \in \mathcal{S}_{1,2,4,6}$, then $C_n$ is not $S$-packing critical 
	if and only if $n = 4$ or $n \ge 8$ and $n = 7k + 8m$ for some $k,m \in \mathbb{Z}_{\ge 0}$.
\end{proposition}

\begin{proof}
	Let $S \in \mathcal S_{1,2,4,6}$. For $S \in \mathcal{S}_{1,2,4,6}$, a path $P_n$ requires at least eight vertices
	in order to admit an $S$-packing coloring using four colors. Let $n \ge 8$.
	Using the repeating color pattern $(12131214)^*$ along $P_n$ up to length $n$, we obtain an $S$-packing coloring.
	Therefore, for all $n \ge 8$, $\chi_S(P_n)=4$ and hence $\chi_S(C_n)\ge 4$.
	
	If $n \ge8$, we consider the following colorings of $C_n$.
	The referred patterns start from vertex $v_1$, and after a specified initial
	sequence, the coloring repeats the pattern $12131214$ so that the color of $v_n$ will be $4$.
	
	If $n \equiv 0 \pmod{8}$, we color $C_n$ with $(12131214)^*$.
	
	If $n \equiv 1 \pmod{8}$, $n \notin \{9, 17, 25, 33, 41\}$ we color $C_n$ as $(1213124)^7(12131214)^*$.
	
	If $n \equiv 2 \pmod{8}$, $n \notin \{10, 18, 26, 34\}$ we color $C_n$ as $(1213124)^6(12131214)^*$.
	
	If $n \equiv 3 \pmod{8}$, $n \notin \{11, 19, 27\}$ we color $C_n$ as $(1213124)^5(12131214)^*$.
	
	If $n \equiv 4 \pmod{8}$, $n \notin \{12, 20\}$ we color $C_n$ as $(1213124)^4(12131214)^*$.
	
	If $n \equiv 5 \pmod{8}$, $n \neq 13 $ we color $C_n$ as  $(1213124)^3(12131214)^*$.
	
	If $n \equiv 6 \pmod{8}$, we color $C_n$ as $(1213124)^2(12131214)^*$.
	
	If $n \equiv 7 \pmod{8}$, we color $C_n$ as $(1213124)(12131214)^*$
	
	It follows that $\chi_S(C_n) \le 4$ and, in turn, $\chi_S(C_n)=4$ for every $n \ge 8, n \notin \{9, 10, 11, 12, 13, \allowbreak 17, 18, 19, 20, 25, 26, 27, 33, 34, 41\}$.
	We conclude that in this case there is no $S$-packing critical cycle.
	
	For the small cases and $n \in \{9, 10, 11, 12, 13, \allowbreak 17, 18, 19, 20, 25, 26, 27, 33, 34, 41\}$, we observe that
	$\chi_S(P_3)=2 < \chi_S(C_3)=3$,
	$\chi_S(P_4)=3=\chi_S(C_4)$,
	$\chi_S(P_5)=3 < \chi_S(C_5)=4$,
	$\chi_S(P_6)=3 < \chi_S(C_6)=4$,
	$\chi_S(P_7)=3 < \chi_S(C_7)=4$,
	$\chi_S(P_8)=4=\chi_S(C_8)$,
	$\chi_S(P_9)=4 < \chi_S(C_9)=5$,
	$\chi_S(P_{10})=4 < \chi_S(C_{10})=5$,
	$\chi_S(P_{11})=4 < \chi_S(C_{11})=5$,
	$\chi_S(P_{12})=4 < \chi_S(C_{12})=5$,
	$\chi_S(P_{13})=4 < \chi_S(C_{13})=5$,
	$\chi_S(P_{17})=4 < \chi_S(C_{17})=5$,
	$\chi_S(P_{18})=4 < \chi_S(C_{18})=5$,
	$\chi_S(P_{19})=4 < \chi_S(C_{19})=5$,
	$\chi_S(P_{20})=4 < \chi_S(C_{20})=5$,
	$\chi_S(P_{25})=4 < \chi_S(C_{25})=5$,
	$\chi_S(P_{26})=4 < \chi_S(C_{26})=5$,
	$\chi_S(P_{27})=4 < \chi_S(C_{27})=5$,
	$\chi_S(P_{33})=4 < \chi_S(C_{33})=5$,
	$\chi_S(P_{34})=4 < \chi_S(C_{34})=5$,
	$\chi_S(P_{41})=4 < \chi_S(C_{41})=5$.
	
	Now, in this case we may conclude that $C_3, C_5, C_6, C_7, C_9, C_{10}, C_{11}, C_{12}, C_{13}, C_{17}, C_{18}, \allowbreak C_{19}, C_{20}, C_{25}, C_{26}, C_{27}, C_{33}, C_{34}$,  and $C_{41}$ are the only $S$-packing critical cycles if $S \in \mathcal S_{1,2,4,6}$.
\end{proof}

\end{document}
